\documentclass[12pt]{article}
\usepackage[utf8]{inputenc}
\usepackage[T1]{fontenc}
\usepackage[margin=1in]{geometry}
\usepackage{setspace}
\usepackage{xcolor}
\usepackage[colorlinks=true, urlcolor=blue!60!black, linkcolor=blue!60!black]{hyperref}
\setlength{\parindent}{0em}
\setlength{\parskip}{0.45em}
\renewcommand{\baselinestretch}{1.0}
\begin{document}
\begin{center}
{\Large \textit{Fitzpatrick v. Bitzer}, 427 U.S. 445 (1976)\par}
\vspace{0.3em}
{\normalsize Case Brief --- One-Page Edition\par}
\end{center}

\textbf{Memory Jogger.} Section 5 of the Fourteenth Amendment lets Congress abrogate state sovereign immunity (\href{https://supreme.justia.com/cases/federal/us/427/445/}{427 U.S. 445}).

\textbf{Facts.} Connecticut let women with 25 years of service retire on full pensions at 50 but made men wait until 55. Male employees sued under Title VII, which the 1972 Amendments, enacted under Section 5, had extended to state employers with backpay and fee remedies.

\textbf{Procedural History.} The district court enjoined the plan but denied backpay and fees under the Eleventh Amendment; the Second Circuit allowed fees only.

\textbf{Votes.} 9--0 in the judgment; Rehnquist wrote for seven, and Brennan and Stevens concurred in the judgment.

\textbf{Holding.} When Congress acts under Section 5 and clearly authorizes suit, federal courts may award retroactive backpay and attorney's fees against States from the state treasury.

\textbf{Reasoning.} The Fourteenth Amendment, ratified after the Eleventh, aims its prohibitions at the States and grants Congress enforcement power, a ``carving out'' of state sovereignty under \textit{Ex parte Virginia}. Brennan would have found no immunity in suits by a State's own citizens; Stevens would have relied on the Commerce Clause. The Article I contrast foreshadowed \textit{Seminole Tribe}.

\textbf{Examples.} \emph{Same-Side:} FMLA family-care damages (\textit{Hibbs}). \emph{Opposite-Side:} Commerce Clause abrogation (\textit{Seminole Tribe}). \emph{Fence-Sitter:} disparate-impact damages against state universities.

\textbf{Critique.} Originalists praise the last-in-time logic but note the text's silence on damages suits; progressives say the later Article I and \textit{City of Boerne} limits turned \textit{Fitzpatrick} into a bottleneck.

\textbf{Key Quotation.} ``But we think that the Eleventh Amendment, and the principle of state sovereignty which it embodies, see \textit{Hans v.\ Louisiana}, 134 U.S.\ 1 (1890), are necessarily limited by the enforcement provisions of \S~5 of the Fourteenth Amendment.'' --- Justice Rehnquist, Opinion of the Court, 427 U.S.\ at 456.

\end{document}
